%%%%%%%%%%%%%%%%
%%% exod 2 %%%
%%%%%%%%%%%%%%%%

\Chapter{2}

\Verse{1}
{
	\hJ{וַיֵּלֶךְ אִישׁ מִבֵּית לֵוִי וַיִּקַּח אֶת בַּת לֵוִי׃}
}
{
	\eJ{
		And a man from the house of Levi went and took a
		daughter of Levi.
	}
}
{
	Two Levites become attached to one another.
	
	\aB{Deeply.}
}

\Verse{2}
{
	\hJ{וַתַּהַר הָאִשָּׁה וַתֵּלֶד בֵּן וַתֵּרֶא אֹתוֹ כִּי טוֹב הוּא וַתִּצְפְּנֵהוּ שְׁלֹשָׁה יְרָחִים׃}
}
{
	\eJ{
		And the woman became pregnant and gave birth to a son. And
		she saw him, that he was good, and she concealed him for
		three months.
	}
}
{
	It's the bible, so immediately she's pregnant.
	
	\aB{All verbs in the bible get you pregnant.}

	It's the illegal sex.\fB{You know the one.}

	But the mom thinks he's great, so she hides him.

%	\footnote{
%		gave birth to a son. Moses is profoundly a lone figure.
%		Although he has
%		a family, Exodus is not about family relations and does
%		not develop them
%		in Moses’ case. Coming on the heels of Genesis, with its
%		long stories of
%		families, this is striking. We learn little of Moses’
%		mother and less of
%		his father. The most central family relationship is
%		between Moses and
%		his brother Aaron, yet it plays no role in the story.
%		Aaron need not be
%		Moses’ brother for the sake of the development of the
%		story; their
%		interactions usually do not depend on it at all. And
%		Miriam, when she is
%		first identified by name, is identified as “the sister
%		of Aaron,” rather
%		than of Moses (15:20), and she and Moses are never
%		pictured exchanging
%		any words. Family members play a part in the birth story
%		of Moses, but
%		the account establishes, after all, not a relationship
%		but the distance
%		between Moses and his family; it is about his being
%		raised by others,
%		from another people. Moses has a family of which he is
%		husband and
%		father, but this just further demonstrates the point,
%		because the text
%		merely reports that he has a wife and sons. They play no
%		role. There is
%		no story about them except the strange story of
%		Zipporah’s circumcising
%		their son, and it is only three verses long (4:24–26),
%		and it is
%		incomprehensible. (See the comment on 4:24.)
%	}
}

\Verse{3}
{
	\hJ{וְלֹא יָכְלָה עוֹד הַצְּפִינוֹ וַתִּקַּח לוֹ תֵּבַת גֹּמֶא וַתַּחְמְרָה בַחֵמָר וּבַזָּפֶת וַתָּשֶׂם בָּהּ אֶת הַיֶּלֶד וַתָּשֶׂם בַּסּוּף עַל שְׂפַת הַיְאֹר׃}
}
{
	\eJ{
		And she was not able to conceal him anymore, and she took
		an ark made of bulrushes for him and smeared it with
		bitumen and with pitch and put the boy in it and put it in
		the reeds by the bank of the Nile.
	}
}
{
	However babies are loud and they get bigger over time so
	hiding them works less well as they age.

	So eventually the mom makes a box%
	\fB{An ark is a box.}%
	\fC{Same Heb. word that was used for Noah's ark in Gen. 6-9.}
	out of this hard grass type stuff and glues it together
	with this dark glue type stuff and this other dark glue
	type stuff and puts the kid in it and puts it in this
	other hard grass type stuff%
	\fC{
		The Heb. does not say she put him ``in the Nile,''
		but \textsl{``in the reeds at the edge of the Nile.''}
		Verse 5 likewise finds the ark ``among the reeds.''
		It may have been in shallow water, but the text never
		explicitly says it was floating down the river.
		This detail has been surprisingly little discussed.
	}
	over by the ground right next to the water.


%	\footnote{
%		the Nile. The same river that means death for all the
%		other male
%		newborns means life for Moses.
%	}
}

\Verse{4}
{
	\hJ{וַתֵּתַצַּב אֲחֹתוֹ מֵרָחֹק לְדֵעָה מַה יֵּעָשֶׂה לוֹ׃}
}
{
	\eJ{
		And his sister stood still at a distance to know what
		would be done to him.
	}
}
{
	His sister stands \aB{still} at a distance to see what happens.%
	\fC{
		Heb. \heb{ותתצב}, from \heb{יצב}, means to station oneself
		or take one's stand. In the bible this verb is never used
		for walking, running, or otherwise moving around;
		it consistently describes taking or holding a stationary
		position. So ``stood still'' captures the sense well.
		This detail has been surprisingly little discussed.
	}

	\aB{
		Many popular depictions of this scene depict the sister
		running after the basket, while the others show her
		walking after it while it drifts away. To our knowledge,
		no depiction of this scene has yet had the
		b\eR{iblic}al\eR{ }l\eR{iterali}s\eR{m or balls} to show
		the scene as it is depicted in the text.

		This little detail is important because it determines
		whether or not the baby in question already lived at
		the pharaoh's house to begin with.
		
		Just make a note of that and let's keep moving.
	}
}

\Verse{5}
{
	\hJ{וַתֵּרֶד בַּת פַּרְעֹה לִרְחֹץ עַל הַיְאֹר וְנַעֲרֹתֶיהָ הֹלְכֹת עַל יַד הַיְאֹר וַתֵּרֶא אֶת הַתֵּבָה בְּתוֹךְ הַסּוּף וַתִּשְׁלַח אֶת אֲמָתָהּ וַתִּקָּחֶהָ׃}
}
{
	\eJ{
		And the Pharaoh’s daughter went down to bathe at the Nile,
		and her girls were going alongside the Nile, and she saw
		the ark among the reeds and sent her maid, and she took
		it.
	}
}
{
	So the king's daughter goes down to the river to take a bath.

	Her girls are goofing around nearby.

	Now the king's daughter sees the basket \aB{among the reeds}\fC{Where it started.}
	and sends the help to fetch it.
}

\Verse{6}
{
	\hJ{וַתִּפְתַּח וַתִּרְאֵהוּ אֶת הַיֶּלֶד וְהִנֵּה נַעַר בֹּכֶה וַתַּחְמֹל עָלָיו וַתֹּאמֶר מִיַּלְדֵי הָעִבְרִים זֶה׃}
}
{
	\eJ{
		And she opened it and saw him, the child: and here was a
		boy crying, and she had compassion on him, and she said,
		“This is one of the Hebrews’ children.”
	}
}
{
	She sees the baby and realizes it's one of the Heeb--- He babies.
}

\Verse{7}
{
	\hJ{וַתֹּאמֶר אֲחֹתוֹ אֶל בַּת פַּרְעֹה הַאֵלֵךְ וְקָרָאתִי לָךְ אִשָּׁה מֵינֶקֶת מִן הָעִבְרִיֹּת וְתֵינִק לָךְ אֶת הַיָּלֶד׃}
}
{
	\eJ{
		And his sister said to Pharaoh’s daughter, “Shall I go and
		call a nursing woman from the Hebrews for you, and she’ll
		nurse the child for you?”
	}
}
{
	Now the sister of the baby is right there where the king's daughter is.
	
	\aB{She got to the palace by standing as fast as she could.}

	So she (the sister) says to the king's daughter,

	Hey, want a free lactating lady?
	
	... I know an authentic Heeb--- Shebrew lady.
}

\Verse{8}
{
	\hJ{וַתֹּאמֶר לָהּ בַּת פַּרְעֹה לֵכִי וַתֵּלֶךְ הָעַלְמָה וַתִּקְרָא אֶת אֵם הַיָּלֶד׃}
}
{
	\eJ{
		And Pharaoh’s daughter said to her, “Go.” And the girl
		went and called the child’s mother.
	}
}
{
	The king's daughter's like ``Ok.''

	The sister calls mom.
}

\Verse{9}
{
	\hJ{וַתֹּאמֶר לָהּ בַּת פַּרְעֹה הֵילִיכִי אֶת הַיֶּלֶד הַזֶּה וְהֵינִקִהוּ לִי וַאֲנִי אֶתֵּן אֶת שְׂכָרֵךְ וַתִּקַּח הָאִשָּׁה הַיֶּלֶד וַתְּנִיקֵהוּ׃}
}
{
	\eJ{
		And Pharaoh’s daughter said to her, “Take this child and
		nurse him for me, and I’ll give your pay.” And the woman
		took the boy and nursed him.
	}
}
{
	So now the king's daughter hires the baby's mom to
	raise her own baby and pays her for it.

	\aB{God I love this book.}

	\aC{
		In summary, contrary to the commonly pictured version
		of this scene, it would appear that the text is hinting,
		in no uncertain terms, that (1) the baby's family already
		lives either at the Pharaoh's palace, or sufficiently close
		that one can walk there without moving, and (2) further,
		it sounds a bit as if the king has been tricked into
		subsidizing this particular Hebrew child's family.
	}
}

\Verse{10}
{
	\hJ{וַיִּגְדַּל הַיֶּלֶד וַתְּבִאֵהוּ לְבַת פַּרְעֹה וַיְהִי לָהּ לְבֵן וַתִּקְרָא שְׁמוֹ מֹשֶׁה וַתֹּאמֶר כִּי מִן הַמַּיִם מְשִׁיתִהוּ׃}
}
{
	\eJ{
		And the boy grew older, and she brought him to Pharaoh’
		daughter, and he became her son. And she called his name
		Moses, and she said, “Because I drew him from the water.”
	}
}
{
	The baby gets older.

	Eventually, the mom brings the baby to Pharaoh's daughter.

	At this point they finally give the baby a name.\fB{
		He was just ``baby'' before this.
	}

	The king's daughter chooses the name, which turns out
	to be (1) an Egyptian throne name common around
	the 1200s B.C., and (2) not described as an Egyptian
	name by the text, even though the Egyptian girl names
	the boy, and (3) described by the text as \textsl{(surprise)}
	a Hebrew pun, and the king's daughter decides
	to name him \textsl{(drumroll)}

	\begin{quote}
	\aB{Drew.}

	\aA{Because I drew him from the water.}

	\aC{
		Heb. \paleo{מְשִׁה} (\textsl{msh}) because \paleo{מְשִׁיתִֽהוּ} (\textsl{mshitho}) from the water.
	}
	\end{quote}

%	\footnote{
%		Pharaoh’s daughter, and he became her son. The story is
%		intriguingly
%		similar to the legend of the Mesopotamian king Sargon,
%		found in Assyrian
%		and Babylonian texts, in which a priestess places her
%		infant son in a
%		basket of rushes with pitch on its exterior and casts
%		him into the
%		river, and a water-drawer retrieves the baby and rears
%		him as his son.
%		This and other literary parallels to the birth account
%		of Moses suggest
%		how enigmatic the biblical story is. Freud observed that
%		such stories
%		generally involve three steps: (1) a child is born of
%		noble or royal
%		lineage, (2) the child comes to be brought up as a
%		commoner, and (3) the
%		child grows up and eventually arrives back at his
%		rightful place in a
%		royal house. Freud noted that such stories were
%		conceived etiologically,
%		composed as justifications of cases in which commoners
%		rose to thrones.
%		The historical truth in such cases lay in the second
%		step: the king
%		really came from commoner roots. The story was composed
%		to legitimize
%		his kingship, as an answer to those who would deny his
%		royal blood.
%		Freud considered the birth story of Moses against this
%		background and
%		suggested the possibility that here, too, the truth
%		behind the etiology
%		lay in the second step, that Moses was an Egyptian, and
%		that the birth
%		story was composed to explain how an Egyptian had come
%		to be the leader
%		of the Israelites. Freud’s interest (on this particular
%		point of his
%		larger study) was primarily historical, and his
%		hypothesis has never
%		been proved or disproved; but it is also important
%		because it indicates
%		what the Torah and its audience valued. In the Israelite
%		story, the
%		values are reversed. The royal house is step two, the
%		aberration, rather
%		than the prized position. Moses’ royal placement simply
%		is not what is
%		important. There is no information about his early life
%		in the Egyptian
%		court. We are informed that he is nursed by his own
%		mother, but we are
%		not certain what this report is supposed to establish.
%		We cannot even be
%		certain that it means that Moses thus knows that he is
%		Israelite. Later
%		the text says that he “went out to his brothers and saw
%		their burdens”
%		(2:11), but this wording, too, is not definitive as to
%		whether he knows
%		that the Israelites are his kin. There is even ambiguity
%		in his killing
%		of an Egyptian taskmaster whom he sees striking a slave:
%		is it because
%		of his bond with the Israelites or his sense of justice?
%		In a curious
%		parallel to the deity Himself, Moses’ background and
%		motives are
%		mysterious. Footnote 2:10. she called his name Moses.
%		Even learned
%		people often recall this story out of order, imagining
%		the Pharaoh’s
%		daughter naming Moses as she draws him from the water,
%		but she does not
%		in fact name him until after he has grown and his mother
%		brings him to
%		her. What was he called until that time? Classical and
%		recent
%		commentators have not addressed this. Since the text
%		does not tell us,
%		we must assume that the concern is to explain the origin
%		of the name
%		Moses and that there is no interest in pursuing whether
%		there was any
%		prior name given by his parents. Presumably the only
%		name known in
%		tradition and history was Moses, and so the author was
%		not free to make
%		up any other. And the story had to ascribe the naming to
%		the Pharaoh’s
%		daughter because she was the one in power, the one who
%		would present him
%		to the Egyptian society, and so on. Still, we can hardly
%		resist
%		wondering why the parents would not be pictured as
%		giving their son a
%		name. If he is with them until he is weaned, that may be
%		several years,
%		and we can hardly resist imagining what the parents
%		would call him. I
%		imagine—this is my midrash—that they do not give him a
%		name. They call
%		him “the child” (Hebrew hayyeled). When they talk to him
%		they call him
%		“my child” (yaldî). And this gives those years a
%		mysterious, portentous
%		quality. They know that his naming simply is not in
%		their power. And so
%		his fate is not in their hands either. Count how many
%		times you call
%		your child by his or her name in a day, and you will
%		know how many times
%		these parents are reminded of their unique situation:
%		thankful that
%		their son alone is spared from death but sad and
%		frightened that he will
%		be raised by others, from the very household that is
%		their enemy, and
%		worried that he will not know who his real people and
%		family are. From
%		the perspective of Jewish history, this is not a
%		singular experience. In
%		the twentieth century Jewish parents in Europe gave
%		their children to
%		non-Jewish families to save them during the holocaust.
%		Footnote 2:10.
%		Because I drew him. The naming of Moses argues both for
%		and against the
%		story’s historicity. On one hand, there is the
%		unlikelihood of the idea
%		that the princess would know Hebrew, let alone choose to
%		derive the
%		baby’s name from a Hebrew etymology. And in fact the
%		name Moses is not
%		Hebrew. It is Egyptian, meaning “is born,” as in the
%		name Ramesses,
%		meaning “Ra (the sun god) bore him.” On the other hand,
%		the fact that
%		the great leader of the Israelites has an Egyptian name
%		(as do other
%		early priests: Phinehas, Hophni) is evidence that
%		Israelites did indeed
%		live for some time in Egypt. Names are valuable evidence
%		in tracing a
%		community’s origins and history. One might suggest that
%		the Egyptian
%		name Moses was just made up to make the story sound
%		authentic. But we
%		can be fairly certain that the Israelites did not make
%		it up, precisely
%		because they told the story about the princess calling
%		him Moses
%		“because I drew him”—which shows that the Israelites
%		were not conscious
%		of the name’s Egyptian meaning!
%	}
}

\Verse{11}
{
	\hJ{וַיְהִי בַּיָּמִים הָהֵם וַיִּגְדַּל מֹשֶׁה וַיֵּצֵא אֶל אֶחָיו וַיַּרְא בְּסִבְלֹתָם וַיַּרְא אִישׁ מִצְרִי מַכֶּה אִישׁ עִבְרִי מֵאֶחָיו׃}
}
{
	\eJ{
		And it was in those days, and Moses grew older, and he
		went out to his brothers and saw their burdens, and he saw
		an Egyptian man striking a Hebrew man, one of his
		brothers.
	}
}
{
	Drew grew.

	Eventually he goes outside and sees the burdens of his
	Hebrew brothers.\fB{Lit. He-bros in the original text.}

	And he sees an Egyptian hitting one of his Hebrew brothers.\fB{Ibid.}

	%	\footnote{
%		Moses. Although Exodus is ultimately about God, Israel,
%		and Egypt, the
%		narrative attention is focused on Moses from its second
%		chapter to the
%		end. Thus, although Exodus is about nations and
%		extraordinary divine
%		interventions into the course of human affairs, it
%		directs its readers
%		to this dynamic through the lens of an individual man.
%		Notably, even
%		though Exodus is not about individuals in the way that
%		Genesis is, it
%		introduces a figure in whom character development
%		reaches a new level,
%		equaled by no other figure in the Hebrew Bible except
%		possibly David.
%		Moses is pictured at various stages in his life,
%		expressing a variety of
%		moods and emotions, changing, especially in the way in
%		which he relates
%		and speaks to God.
%	}
}

\Verse{12}
{
	\hJ{וַיִּפֶן כֹּה וָכֹה וַיַּרְא כִּי אֵין אִישׁ וַיַּךְ אֶת הַמִּצְרִי וַיִּטְמְנֵהוּ בַּחוֹל׃}
}
{
	\eJ{
		And he turned this way and that way and saw that there was
		no man, and he struck the Egyptian and hid him in the
		sand.
	}
}
{
	Drew looks around to make sure no one's watching
	and he kills the Egyptian and hides the dead body
	in the sand.\fB{They have sand in Egypt.}
}

\Verse{13}
{
	\hJ{וַיֵּצֵא בַּיּוֹם הַשֵּׁנִי וְהִנֵּה שְׁנֵי אֲנָשִׁים עִבְרִים נִצִּים וַיֹּאמֶר לָרָשָׁע לָמָּה תַכֶּה רֵעֶךָ׃}
}
{
	\eJ{
		And he went out on the second day, and here were two
		Hebrew men fighting. And he said to the one who was in the
		wrong, “Why do you strike your companion?”
	}
}
{
	A couple days later, he sees two of his bros\fB{Ibid if you skip the sand one.} there fighting.

	He tries to break up the fight.
}

\Verse{14}
{
	\hJ{וַיֹּאמֶר מִי שָׂמְךָ לְאִישׁ שַׂר וְשֹׁפֵט עָלֵינוּ הַלְהרְגֵנִי אַתָּה אֹמֵר כַּאֲשֶׁר הָרַגְתָּ אֶת הַמִּצְרִי וַיִּירָא מֹשֶׁה וַיֹּאמַר אָכֵן נוֹדַע הַדָּבָר׃}
}
{
	\eJ{
		And he said, “Who made you a commander and judge over us?
		Are you saying you’d kill me—the way you killed the
		Egyptian?!” And Moses was afraid and said, “The thing is
		known for sure.”
	}
}
{
	And one of the bros there replies,

	What are you gonna kill us like you killed that other guy?

	\begin{center}
	\eRB{
		And just then, the camera that's filming the bible
		
		closely drew in on Drew and Drew knew
		
		they were onto him, and so Drew
		
		spoke unto the camera saying,

		... Dudetheyreontome.
	}%
	\fC{%
		This line appears to be a late scribal addition,
		likely in service of some theological agenda
		the finer points of which have since been lost
		to history.
	}
	\end{center}%

}

\Verse{15}
{
	\hJ{וַיִּשְׁמַע פַּרְעֹה אֶת הַדָּבָר הַזֶּה וַיְבַקֵּשׁ לַהֲרֹג אֶת מֹשֶׁה וַיִּבְרַח מֹשֶׁה מִפְּנֵי פַרְעֹה וַיֵּשֶׁב בְּאֶרֶץ מִדְיָן וַיֵּשֶׁב עַל הַבְּאֵר׃}
}
{
	\eJ{
		And Pharaoh heard this thing and sought to kill Moses, and
		Moses fled from Pharaoh’s presence and lived in the land
		of Midian. And he sat by a well.
	}
}
{
	And indeed, they were onto him.

	And he knew they were onto him.

	He was thereafter On The Lamb.\fC{Sic.}

	\aB{
		And Drew said ``oh well'' for the king sought
		his head so with heart full of dred he drew
		himšelf out of bed and he fled about five
		hundred miles to Midian to avoid being dead
		and when he gets there he sits by a well as
		a pun.
	}
}

\Verse{16}
{
	\hJ{וּלְכֹהֵן מִדְיָן שֶׁבַע בָּנוֹת וַתָּבֹאנָה וַתִּדְלֶנָה וַתְּמַלֶּאנָה אֶת הָרְהָטִים לְהַשְׁקוֹת צֹאן אֲבִיהֶן׃}
}
{
	\eJ{
		And a priest of Midian had seven daughters, and they came
		and drew water and filled the troughs to water their
		father’s flock,
	}
}
{
	\aB{
		The attentive reader will recall from Genesis that
		wells in the bible are \eR{a} deep \eR{source of water and}
		wet\eR{ness imagery and are deeply associated with
		life, birth,} and if you go to them you immediately
		end up knee deep in p\eRJE{rimary spots for matrimon}y.
	}

	Seven girls come to the hole of wetness
	and a series of drew-related puns happen.
}

\Verse{17}
{
	\hJ{וַיָּבֹאוּ הָרֹעִים וַיְגָרְשׁוּם וַיָּקם מֹשֶׁה וַיּוֹשִׁעָן וַיַּשְׁקְ אֶת צֹאנָם׃}
}
{
	\eJ{
		and the shepherds came and drove them away. And Moses got
		up and saved them and watered their flock.
	}
}
{
	\aB{Here J tells a \eR{w}hole story in like half a sentence.}

	Ok so apparently while the seven girls water was in the process
	of being drewn, some sheep herds heard that the girls were
	getting \eR{water} down there and they came to them
	and chased them away for some reason, nobody knows
	why, this is mostly exegesis and admittedly a
	stretch.

	\aB{Drew saves them.}
}

\Verse{18}
{
	\hJ{וַתָּבֹאנָה אֶל רְעוּאֵל אֲבִיהֶן וַיֹּאמֶר מַדּוּעַ מִהַרְתֶּן בֹּא הַיּוֹם׃}
}
{
	\eJ{
		And they came to Reuel, their father, and he said, “Why
		were you so quick to come today?”
	}
}
{
	So they go back home to their dad.

	\aB{His name is Reuel everyone. Don't forget Reuel.}

	He's like ``That was quick. Why home so early?''
}

\Verse{19}
{
	\hJ{וַתֹּאמַרְןָ אִישׁ מִצְרִי הִצִּילָנוּ מִיַּד הָרֹעִים וְגַם דָּלֹה דָלָה לָנוּ וַיַּשְׁקְ אֶת הַצֹּאן׃}
}
{
	\eJ{
		And they said, “An Egyptian man rescued us from the
		shepherds’ hand, and he drew water for us and watered the
		flock, too.”
	}
}
{
	And the seven girls reply all in unison, something like
	
	Ok so \aB{Drew drew} \aC{water from the} well
	no first there were these men \aB{an Egyptian}\fB{Drew's Egyptian.}
	no \textsl{men} \aB{who Drew} \aC{no the shepherds} and Drew
	\aB{drove them away} because first before Drew
	they were chasing us away from the \aB{well but well Drew}
	\aB{chased them away too} from us seven so they couldn't
	chase us away anymore \aC{and but so} Drew drew water
	\aC{for us too} no us seven \aC{and Drew drew} \aB{water}
	\aC{you talking about} \aB{him} and \aB{t}he
	flock a\aB{n}\aC{d} us too.%
	\fC{The text is a bit much here.}

	There was a lot going on in their explanation but dad's used to it.
}

\Verse{20}
{
	\hJ{וַיֹּאמֶר אֶל בְּנֹתָיו וְאַיּוֹ לָמָּה זֶּה עֲזַבְתֶּן אֶת הָאִישׁ קִרְאֶן לוֹ וְיֹאכַל לָחֶם׃}
}
{
	\eJ{
		And he said to his daughters, “And where is he? Why is
		this that you’ve left the man? Call him, and let him eat
		bread.”
	}
}
{
	So the dad's like ``...You just left him?

	\aB{Now I imagine the girls suddenly pause all like {\reubenface}}

	Have him over for dinner for God's sake.''
}

\Verse{21}
{
	\hJ{וַיּוֹאֶל מֹשֶׁה לָשֶׁבֶת אֶת הָאִישׁ וַיִּתֵּן אֶת צִפֹּרָה בִתּוֹ לְמֹשֶׁה׃}
}
{
	\eJ{
		And Moses was content to live with the man. And he gave
		Zipporah, his daughter, to Moses,
	}
}
{
	Drew comes and \eR{h}e\eR{ e}ats dinner and \eR{marries} one of the girls too.
}

\Verse{22}
{
	\hJ{וַתֵּלֶד בֵּן וַיִּקְרָא אֶת שְׁמוֹ גֵּרְשֹׁם כִּי אָמַר גֵּר הָיִיתִי בְּאֶרֶץ נכְרִיָּה׃}
}
{
	\eJ{
		and she gave birth to a son, and he called his name
		Gershom, “because,” he said, “I was an alien in a foreign
		land.”
	}
}
{
	She's already pre--- \aB{It's a boy!}\fA{The text is a bit rushed here.}

	And this is the bible so they name him a pun.

	Drew names him Foreigner, fore ignerant was he of the customs of that land.
}

\Verse{23}
{
	\hJ{וַיְהִי בַיָּמִים הָרַבִּים הָהֵם וַיָּמת מֶלֶךְ מִצְרַיִם}
	\hP{וַיֵּאָנְחוּ בְנֵי יִשְׂרָאֵל מִן הָעֲבֹדָה וַיִּזְעָקוּ וַתַּעַל שַׁוְעָתָם אֶל הָאֱלֹהִים מִן הָעֲבֹדָה׃}
}
{
	\eJ{And it was after those many days, and the king of Egypt
		died.}
	\eP{And the children of Israel groaned from the work,
		and they cried out, and their wail went up to God from the
		work.}
}
{

	Anyways after all that, the new boss back home dies.

	\aB{So it sounds like things might take a turn for the better and%
	\fC{Lit \texttt{\textbar} in the original, a vertical stroke which can represent
	either ``and'' in Hebrew or ``or'' in the orthographies said to
	be used among people of ``Heshbon and Heshbon'' (Heb. \heb{חֶשְׁבּוֹןוחֶשְׁבּוֹן}).
	The Septuagint renders the names of these lands as
	Logos and Logismos (Grk. λόγος καὶ λογισμός).
	The Vulgate has Ratiō et Computātiō, all generally understood
	to refer to the same place.
	The location of biblical Heshbon has not been identified.
	}
	maybe not}

	\aC{
		So last we left off back in Ex. 1:14 we heard Egypt
		was making the children of Israel's life
	}
	\aB{suck}
	\aC{
		hard with work and with bricks and also with the glue that
		goes between the bricks—
	}
	\aB{however}
	\eR{they could do so}
	\aC{they did to them, they did so}
	\aC{—with all the work in the
		field—all their work that they did for them—with harshness.
	}
	\eR{But }
	\aB{
		in\eR{ their }stead\eR{fast faith they held fast}
		of\eR{ten onto their faith in} that
		\eR{land. This completes the recap of where we left off,
		and having provided that} back\eR{ground, you're
		in luck, for we now turn} to P.
	}
}

\Verse{24}
{
	\hP{וַיִּשְׁמַע אֱלֹהִים אֶת נַאֲקָתָם וַיִּזְכֹּר אֱלֹהִים אֶת בְּרִיתוֹ אֶת אַבְרָהָם אֶת יִצְחָק וְאֶת יַעֲקֹב׃}
}
{
	\eP{
		And God heard their moaning, and God remembered His
		covenant with Abraham, with Isaac, and with Jacob.
	}
}
{
%	\footnote{
%		groan, cry, wail, moan. Four different words are used in
%		the Hebrew to
%		describe their crying. This conveys that their agony is
%		intense,
%		continuous, and pervasive.
%	}
}

\Verse{25}
{
	\hP{וַיַּרְא אֱלֹהִים אֶת בְּנֵי יִשְׂרָאֵל וַיֵּדַע אֱלֹהִים׃}
}
{
	\eP{And God saw the children of Israel. And God knew!}
}
{
}
